\documentclass[a4paper, 12pt]{article}
\usepackage[T2A]{fontenc}
\usepackage[english,russian]{babel}
\usepackage{amsmath}
\usepackage{amsthm}
\usepackage{amsfonts}
\usepackage{amssymb}
\usepackage{pgfplots}


% Абзацы и геометрия страницы
\usepackage{indentfirst}
%\usepackage{geometry}
\usepackage[
    top=2cm,
    bottom=2cm,
    left=1cm,
    right=1cm
]{geometry}


\begin{document}

\begin{center}
    МИНИСТЕРСТВО НАУКИ И ВЫСШЕГО ОБРАЗОВАНИЯ РОССИЙСКОЙ ФЕДЕРАЦИИ

    ГОСУДАРСТВЕННОЕ ОБРАЗОВАТЕЛЬНОЕ УЧРЕЖДЕНИЕ ОБРАЗОВАНИЯ

    РЯЗАНСКИЙ ГОСУДАРСТВЕННЫЙ РАДИОТЕХНИЧЕСКИЙ УНИВЕРСИТЕТ 

    ИМ. В. Ф. УТКИНА

    КАФЕДРА ВЫЧИСЛИТЕЛЬНОЙ И ПРИКЛАДНОЙ МАТЕМАТИКИ 
\end{center}

\vspace{5 cm}

\begin{center}
    ЛАБОРАТОРНАЯ РАБОТА № 2

    ПО ДИСЦИПЛИНЕ "ТЕОРЕТИЧЕСКАЯ ИНФОРМАТИКА"

    НА ТЕМУ "КОЛИЧЕСТВО ИНФОРМАЦИИ. ОБЪЕМНЫЙ ПОДХОД"

    ВАРИАНТ 14

\end{center}

\vspace{7 cm}

\begin{flushright}
    Выполнил:

    Студент группы 6413

    Косяков Д.А.

    Проверил

    Доцент кафедры ВПМ

    Никичкин Борис Владимирович
\end{flushright}

\begin{center}
    РЯЗАНЬ 2026
\end{center}

\pagebreak

\textbf{ Цель работы:} измерить количество информации согласно объемному подходу к измерению информации.

\textbf{ Задание:} Существуют два типа объектов А и В, которые хранятся побитно. N объектов А занимают объем Х байт, а М объектов В --- 
объем Y Кбайт.
Какой суммарный объем занимают К объектов А и L объектов В? Результат представить в байтах, Кбайтах и Мбайтах (округлить до сотых).
Размер объектов А и В в битах округлить в меньшую сторону.

\textbf{ Выполнение задания:} заданы N = 13, X = 31, M = 6718, Y = 18.86, K = 5283, L = 9326. 

Определим размер одного объекта А. Объем Х указан в байтах, поэтому его нужно перевести в биты, умножив на 8: 
$$V_A = \frac{X \cdot 8}{N} = \frac{31 \cdot 8}{13} =  \frac{248}{13} = 19.0769 \approx 19 \text{ бит. }$$

Определив размер одного объекта В. Объем Y указан в Кбайтах, поэтому его нужно перевести в биты, умножив на $8 \cdot 1024 = 8192$. 
Результат округлим в меньшую сторону: $$V_B = \frac{Y \cdot 8192}{M} = \frac{18.86 \cdot 8192}{6718} = \frac{154501.12}{6718} = 22.998 \text{ бит } \approx 22 \text{ бит. }$$

 Определим суммарный объем, который занимают K = 5283 объектов А и L = 9326 объектов В.
 $$V = K \cdot V_{A} + l \cdot V_B = 5283 \cdot 19 + 9326 \cdot 22 = 305549 \text{ бит. }$$
 Переведем объем в байты: $$V = \frac{305549}{8} = 38194 \text{ байт. }$$
 Переведем в Кбайты: $$V = \frac{38194}{1024} = 37.3 \text{ Кбайт. }$$
 Переведем в Мбайты: $$V = \frac{37.30}{1024} = 0.04 \text{ Мбайт. }$$

\end{document}